% =============================================================================
% Shared preamble for the MIPT FPMI master's-entrance prep textbook.
%
% Loaded by both book-ru.tex and book-en.tex. The active language is selected
% by each book file via babel's option list (last language wins) and an
% explicit \selectlanguage{} after % =============================================================================
% Shared preamble for the MIPT FPMI master's-entrance prep textbook.
%
% Loaded by both book-ru.tex and book-en.tex. The active language is selected
% by each book file via babel's option list (last language wins) and an
% explicit \selectlanguage{} after % =============================================================================
% Shared preamble for the MIPT FPMI master's-entrance prep textbook.
%
% Loaded by both book-ru.tex and book-en.tex. The active language is selected
% by each book file via babel's option list (last language wins) and an
% explicit \selectlanguage{} after % =============================================================================
% Shared preamble for the MIPT FPMI master's-entrance prep textbook.
%
% Loaded by both book-ru.tex and book-en.tex. The active language is selected
% by each book file via babel's option list (last language wins) and an
% explicit \selectlanguage{} after \input{preamble}.
%
% Build: pdflatex + biber. See README.md and scripts/build.py.
% Author note: keep this file in English — it's tooling, not content.
% =============================================================================

% --- encoding & languages -----------------------------------------------------
\usepackage[utf8]{inputenc}
\usepackage[T2A,T1]{fontenc}
% Both languages loaded; babel makes the LAST one the main language
% (here Russian) so cyrillic fontencoding bindings biblatex-gost relies
% on are set up at \begin{document}. Each book file then issues
% \selectlanguage{...} immediately inside \begin{document} to set its
% own active language for captions, hyphenation, dates, etc.
\usepackage[english,russian]{babel}

% --- page geometry ------------------------------------------------------------
\usepackage[a4paper, margin=2.5cm]{geometry}

% --- math packages ------------------------------------------------------------
\usepackage{amsmath, amssymb, amsthm, mathtools}
\usepackage{mathrsfs}     % \mathscr
\usepackage{bm}           % bold math

% --- math "house style" macros ------------------------------------------------
% Some of these short names (notably \C from T1 fontenc, and \Prob from
% other packages) are already defined when the preamble runs, so we clear
% them first before redefining. If you need the original meaning of \C
% (caron accent) elsewhere, use \v{c} instead.
\let\R\undefined
\let\C\undefined
\let\N\undefined
\let\Z\undefined
\let\Q\undefined
\let\F\undefined
\let\E\undefined
\let\Prob\undefined

% Number sets.
\newcommand{\R}{\mathbb{R}}
\newcommand{\C}{\mathbb{C}}
\newcommand{\N}{\mathbb{N}}
\newcommand{\Z}{\mathbb{Z}}
\newcommand{\Q}{\mathbb{Q}}
\newcommand{\F}{\mathbb{F}}

% Probability symbols.
\newcommand{\E}{\mathbb{E}}
\newcommand{\Prob}{\mathbb{P}}

% Operators (rendered upright).
\DeclareMathOperator{\rank}{rank}
\DeclareMathOperator{\im}{Im}
\DeclareMathOperator{\spn}{span}
\DeclareMathOperator{\diag}{diag}
\DeclareMathOperator{\tr}{tr}
\DeclareMathOperator{\sgn}{sgn}
\DeclareMathOperator{\Var}{Var}
\DeclareMathOperator{\Cov}{Cov}
\DeclareMathOperator*{\argmin}{arg\,min}
\DeclareMathOperator*{\argmax}{arg\,max}

% Vectors and matrices.
\newcommand{\vect}[1]{\mathbf{#1}}
\newcommand{\mat}[1]{\mathbf{#1}}
\newcommand{\T}{^{\mathsf{T}}}     % transpose

% --- theorem environments -----------------------------------------------------
% Names are localized via babel's \iflanguage in the headers below.
\theoremstyle{plain}
\newtheorem{theorem}{\iflanguage{russian}{Теорема}{Theorem}}[chapter]
\newtheorem{lemma}[theorem]{\iflanguage{russian}{Лемма}{Lemma}}
\newtheorem{proposition}[theorem]{\iflanguage{russian}{Утверждение}{Proposition}}
\newtheorem{corollary}[theorem]{\iflanguage{russian}{Следствие}{Corollary}}

\theoremstyle{definition}
\newtheorem{definition}[theorem]{\iflanguage{russian}{Определение}{Definition}}
\newtheorem{example}[theorem]{\iflanguage{russian}{Пример}{Example}}
\newtheorem{problem}[theorem]{\iflanguage{russian}{Задача}{Problem}}

\theoremstyle{remark}
\newtheorem*{remark}{\iflanguage{russian}{Замечание}{Remark}}

% Localized strings (proofname, cleveref names) are bundled into
% \setupEnglish / \setupRussian below — see end of preamble.

% --- ticket environment -------------------------------------------------------
% Renders the verbatim formulation from программа.pdf at the top of each topic.
% Visual styling only — language-aware label is unnecessary because each topic
% file is per-language already.
\newenvironment{ticket}{%
  \par\medskip\noindent\itshape%
}{%
  \par\medskip%
}

% --- bibliography (GOST 7.0.5/7.1) --------------------------------------------
\usepackage[
    backend=biber,
    style=gost-numeric,
    sorting=none,
    language=auto,
    autolang=other,
]{biblatex}
% Resolved relative to the compilation cwd. The build script (scripts/build.py)
% adds the repo root to BIBINPUTS, so this works for both book-* (cwd = root)
% and chapter-* (cwd = chapter dir, subfile compilation).
\addbibresource{literature.bib}

% --- hyperlinks & cross-refs --------------------------------------------------
\usepackage{hyperref}
\hypersetup{
    colorlinks=true,
    linkcolor=blue!50!black,
    citecolor=green!50!black,
    urlcolor=blue!70!black,
}
\usepackage[capitalize,nameinlink]{cleveref}

% --- Language setup commands -------------------------------------------------
% Each book file calls one of these immediately inside \begin{document}.
% They bundle babel's language switch with proof- and cleveref-name
% localization, so the result is the same regardless of which language
% babel happened to pick as "main".
\newcommand{\setupRussian}{%
  \selectlanguage{russian}%
  \renewcommand{\proofname}{Доказательство}%
  \crefname{theorem}{теорема}{теоремы}\Crefname{theorem}{Теорема}{Теоремы}%
  \crefname{lemma}{лемма}{леммы}\Crefname{lemma}{Лемма}{Леммы}%
  \crefname{proposition}{утверждение}{утверждения}\Crefname{proposition}{Утверждение}{Утверждения}%
  \crefname{corollary}{следствие}{следствия}\Crefname{corollary}{Следствие}{Следствия}%
  \crefname{definition}{определение}{определения}\Crefname{definition}{Определение}{Определения}%
  \crefname{example}{пример}{примеры}\Crefname{example}{Пример}{Примеры}%
  \crefname{problem}{задача}{задачи}\Crefname{problem}{Задача}{Задачи}%
  \crefname{remark}{замечание}{замечания}\Crefname{remark}{Замечание}{Замечания}%
  \crefname{section}{раздел}{разделы}\Crefname{section}{Раздел}{Разделы}%
  \crefname{chapter}{глава}{главы}\Crefname{chapter}{Глава}{Главы}%
  \crefname{equation}{уравнение}{уравнения}\Crefname{equation}{Уравнение}{Уравнения}%
}
\newcommand{\setupEnglish}{%
  \selectlanguage{english}%
  \renewcommand{\proofname}{Proof}%
  \crefname{theorem}{theorem}{theorems}\Crefname{theorem}{Theorem}{Theorems}%
  \crefname{lemma}{lemma}{lemmas}\Crefname{lemma}{Lemma}{Lemmas}%
  \crefname{proposition}{proposition}{propositions}\Crefname{proposition}{Proposition}{Propositions}%
  \crefname{corollary}{corollary}{corollaries}\Crefname{corollary}{Corollary}{Corollaries}%
  \crefname{definition}{definition}{definitions}\Crefname{definition}{Definition}{Definitions}%
  \crefname{example}{example}{examples}\Crefname{example}{Example}{Examples}%
  \crefname{problem}{problem}{problems}\Crefname{problem}{Problem}{Problems}%
  \crefname{remark}{remark}{remarks}\Crefname{remark}{Remark}{Remarks}%
  \crefname{section}{section}{sections}\Crefname{section}{Section}{Sections}%
  \crefname{chapter}{chapter}{chapters}\Crefname{chapter}{Chapter}{Chapters}%
  \crefname{equation}{equation}{equations}\Crefname{equation}{Equation}{Equations}%
}

% --- subfiles: each chapter compiles as a standalone PDF too ------------------
\usepackage{subfiles}

% --- misc ---------------------------------------------------------------------
\usepackage{enumitem}
\usepackage{booktabs}
% Font expansion needs scalable Cyrillic fonts (the `cm-super` Debian package
% or `texlive-fonts-extra`). Disable it to stay portable; protrusion still
% works with the default Computer Modern bitmap fonts.
\usepackage[expansion=false,protrusion=true,final]{microtype}
\usepackage{tikz}
.
%
% Build: pdflatex + biber. See README.md and scripts/build.py.
% Author note: keep this file in English — it's tooling, not content.
% =============================================================================

% --- encoding & languages -----------------------------------------------------
\usepackage[utf8]{inputenc}
\usepackage[T2A,T1]{fontenc}
% Both languages loaded; babel makes the LAST one the main language
% (here Russian) so cyrillic fontencoding bindings biblatex-gost relies
% on are set up at \begin{document}. Each book file then issues
% \selectlanguage{...} immediately inside \begin{document} to set its
% own active language for captions, hyphenation, dates, etc.
\usepackage[english,russian]{babel}

% --- page geometry ------------------------------------------------------------
\usepackage[a4paper, margin=2.5cm]{geometry}

% --- math packages ------------------------------------------------------------
\usepackage{amsmath, amssymb, amsthm, mathtools}
\usepackage{mathrsfs}     % \mathscr
\usepackage{bm}           % bold math

% --- math "house style" macros ------------------------------------------------
% Some of these short names (notably \C from T1 fontenc, and \Prob from
% other packages) are already defined when the preamble runs, so we clear
% them first before redefining. If you need the original meaning of \C
% (caron accent) elsewhere, use \v{c} instead.
\let\R\undefined
\let\C\undefined
\let\N\undefined
\let\Z\undefined
\let\Q\undefined
\let\F\undefined
\let\E\undefined
\let\Prob\undefined

% Number sets.
\newcommand{\R}{\mathbb{R}}
\newcommand{\C}{\mathbb{C}}
\newcommand{\N}{\mathbb{N}}
\newcommand{\Z}{\mathbb{Z}}
\newcommand{\Q}{\mathbb{Q}}
\newcommand{\F}{\mathbb{F}}

% Probability symbols.
\newcommand{\E}{\mathbb{E}}
\newcommand{\Prob}{\mathbb{P}}

% Operators (rendered upright).
\DeclareMathOperator{\rank}{rank}
\DeclareMathOperator{\im}{Im}
\DeclareMathOperator{\spn}{span}
\DeclareMathOperator{\diag}{diag}
\DeclareMathOperator{\tr}{tr}
\DeclareMathOperator{\sgn}{sgn}
\DeclareMathOperator{\Var}{Var}
\DeclareMathOperator{\Cov}{Cov}
\DeclareMathOperator*{\argmin}{arg\,min}
\DeclareMathOperator*{\argmax}{arg\,max}

% Vectors and matrices.
\newcommand{\vect}[1]{\mathbf{#1}}
\newcommand{\mat}[1]{\mathbf{#1}}
\newcommand{\T}{^{\mathsf{T}}}     % transpose

% --- theorem environments -----------------------------------------------------
% Names are localized via babel's \iflanguage in the headers below.
\theoremstyle{plain}
\newtheorem{theorem}{\iflanguage{russian}{Теорема}{Theorem}}[chapter]
\newtheorem{lemma}[theorem]{\iflanguage{russian}{Лемма}{Lemma}}
\newtheorem{proposition}[theorem]{\iflanguage{russian}{Утверждение}{Proposition}}
\newtheorem{corollary}[theorem]{\iflanguage{russian}{Следствие}{Corollary}}

\theoremstyle{definition}
\newtheorem{definition}[theorem]{\iflanguage{russian}{Определение}{Definition}}
\newtheorem{example}[theorem]{\iflanguage{russian}{Пример}{Example}}
\newtheorem{problem}[theorem]{\iflanguage{russian}{Задача}{Problem}}

\theoremstyle{remark}
\newtheorem*{remark}{\iflanguage{russian}{Замечание}{Remark}}

% Localized strings (proofname, cleveref names) are bundled into
% \setupEnglish / \setupRussian below — see end of preamble.

% --- ticket environment -------------------------------------------------------
% Renders the verbatim formulation from программа.pdf at the top of each topic.
% Visual styling only — language-aware label is unnecessary because each topic
% file is per-language already.
\newenvironment{ticket}{%
  \par\medskip\noindent\itshape%
}{%
  \par\medskip%
}

% --- bibliography (GOST 7.0.5/7.1) --------------------------------------------
\usepackage[
    backend=biber,
    style=gost-numeric,
    sorting=none,
    language=auto,
    autolang=other,
]{biblatex}
% Resolved relative to the compilation cwd. The build script (scripts/build.py)
% adds the repo root to BIBINPUTS, so this works for both book-* (cwd = root)
% and chapter-* (cwd = chapter dir, subfile compilation).
\addbibresource{literature.bib}

% --- hyperlinks & cross-refs --------------------------------------------------
\usepackage{hyperref}
\hypersetup{
    colorlinks=true,
    linkcolor=blue!50!black,
    citecolor=green!50!black,
    urlcolor=blue!70!black,
}
\usepackage[capitalize,nameinlink]{cleveref}

% --- Language setup commands -------------------------------------------------
% Each book file calls one of these immediately inside \begin{document}.
% They bundle babel's language switch with proof- and cleveref-name
% localization, so the result is the same regardless of which language
% babel happened to pick as "main".
\newcommand{\setupRussian}{%
  \selectlanguage{russian}%
  \renewcommand{\proofname}{Доказательство}%
  \crefname{theorem}{теорема}{теоремы}\Crefname{theorem}{Теорема}{Теоремы}%
  \crefname{lemma}{лемма}{леммы}\Crefname{lemma}{Лемма}{Леммы}%
  \crefname{proposition}{утверждение}{утверждения}\Crefname{proposition}{Утверждение}{Утверждения}%
  \crefname{corollary}{следствие}{следствия}\Crefname{corollary}{Следствие}{Следствия}%
  \crefname{definition}{определение}{определения}\Crefname{definition}{Определение}{Определения}%
  \crefname{example}{пример}{примеры}\Crefname{example}{Пример}{Примеры}%
  \crefname{problem}{задача}{задачи}\Crefname{problem}{Задача}{Задачи}%
  \crefname{remark}{замечание}{замечания}\Crefname{remark}{Замечание}{Замечания}%
  \crefname{section}{раздел}{разделы}\Crefname{section}{Раздел}{Разделы}%
  \crefname{chapter}{глава}{главы}\Crefname{chapter}{Глава}{Главы}%
  \crefname{equation}{уравнение}{уравнения}\Crefname{equation}{Уравнение}{Уравнения}%
}
\newcommand{\setupEnglish}{%
  \selectlanguage{english}%
  \renewcommand{\proofname}{Proof}%
  \crefname{theorem}{theorem}{theorems}\Crefname{theorem}{Theorem}{Theorems}%
  \crefname{lemma}{lemma}{lemmas}\Crefname{lemma}{Lemma}{Lemmas}%
  \crefname{proposition}{proposition}{propositions}\Crefname{proposition}{Proposition}{Propositions}%
  \crefname{corollary}{corollary}{corollaries}\Crefname{corollary}{Corollary}{Corollaries}%
  \crefname{definition}{definition}{definitions}\Crefname{definition}{Definition}{Definitions}%
  \crefname{example}{example}{examples}\Crefname{example}{Example}{Examples}%
  \crefname{problem}{problem}{problems}\Crefname{problem}{Problem}{Problems}%
  \crefname{remark}{remark}{remarks}\Crefname{remark}{Remark}{Remarks}%
  \crefname{section}{section}{sections}\Crefname{section}{Section}{Sections}%
  \crefname{chapter}{chapter}{chapters}\Crefname{chapter}{Chapter}{Chapters}%
  \crefname{equation}{equation}{equations}\Crefname{equation}{Equation}{Equations}%
}

% --- subfiles: each chapter compiles as a standalone PDF too ------------------
\usepackage{subfiles}

% --- misc ---------------------------------------------------------------------
\usepackage{enumitem}
\usepackage{booktabs}
% Font expansion needs scalable Cyrillic fonts (the `cm-super` Debian package
% or `texlive-fonts-extra`). Disable it to stay portable; protrusion still
% works with the default Computer Modern bitmap fonts.
\usepackage[expansion=false,protrusion=true,final]{microtype}
\usepackage{tikz}
.
%
% Build: pdflatex + biber. See README.md and scripts/build.py.
% Author note: keep this file in English — it's tooling, not content.
% =============================================================================

% --- encoding & languages -----------------------------------------------------
\usepackage[utf8]{inputenc}
\usepackage[T2A,T1]{fontenc}
% Both languages loaded; babel makes the LAST one the main language
% (here Russian) so cyrillic fontencoding bindings biblatex-gost relies
% on are set up at \begin{document}. Each book file then issues
% \selectlanguage{...} immediately inside \begin{document} to set its
% own active language for captions, hyphenation, dates, etc.
\usepackage[english,russian]{babel}

% --- page geometry ------------------------------------------------------------
\usepackage[a4paper, margin=2.5cm]{geometry}

% --- math packages ------------------------------------------------------------
\usepackage{amsmath, amssymb, amsthm, mathtools}
\usepackage{mathrsfs}     % \mathscr
\usepackage{bm}           % bold math

% --- math "house style" macros ------------------------------------------------
% Some of these short names (notably \C from T1 fontenc, and \Prob from
% other packages) are already defined when the preamble runs, so we clear
% them first before redefining. If you need the original meaning of \C
% (caron accent) elsewhere, use \v{c} instead.
\let\R\undefined
\let\C\undefined
\let\N\undefined
\let\Z\undefined
\let\Q\undefined
\let\F\undefined
\let\E\undefined
\let\Prob\undefined

% Number sets.
\newcommand{\R}{\mathbb{R}}
\newcommand{\C}{\mathbb{C}}
\newcommand{\N}{\mathbb{N}}
\newcommand{\Z}{\mathbb{Z}}
\newcommand{\Q}{\mathbb{Q}}
\newcommand{\F}{\mathbb{F}}

% Probability symbols.
\newcommand{\E}{\mathbb{E}}
\newcommand{\Prob}{\mathbb{P}}

% Operators (rendered upright).
\DeclareMathOperator{\rank}{rank}
\DeclareMathOperator{\im}{Im}
\DeclareMathOperator{\spn}{span}
\DeclareMathOperator{\diag}{diag}
\DeclareMathOperator{\tr}{tr}
\DeclareMathOperator{\sgn}{sgn}
\DeclareMathOperator{\Var}{Var}
\DeclareMathOperator{\Cov}{Cov}
\DeclareMathOperator*{\argmin}{arg\,min}
\DeclareMathOperator*{\argmax}{arg\,max}

% Vectors and matrices.
\newcommand{\vect}[1]{\mathbf{#1}}
\newcommand{\mat}[1]{\mathbf{#1}}
\newcommand{\T}{^{\mathsf{T}}}     % transpose

% --- theorem environments -----------------------------------------------------
% Names are localized via babel's \iflanguage in the headers below.
\theoremstyle{plain}
\newtheorem{theorem}{\iflanguage{russian}{Теорема}{Theorem}}[chapter]
\newtheorem{lemma}[theorem]{\iflanguage{russian}{Лемма}{Lemma}}
\newtheorem{proposition}[theorem]{\iflanguage{russian}{Утверждение}{Proposition}}
\newtheorem{corollary}[theorem]{\iflanguage{russian}{Следствие}{Corollary}}

\theoremstyle{definition}
\newtheorem{definition}[theorem]{\iflanguage{russian}{Определение}{Definition}}
\newtheorem{example}[theorem]{\iflanguage{russian}{Пример}{Example}}
\newtheorem{problem}[theorem]{\iflanguage{russian}{Задача}{Problem}}

\theoremstyle{remark}
\newtheorem*{remark}{\iflanguage{russian}{Замечание}{Remark}}

% Localized strings (proofname, cleveref names) are bundled into
% \setupEnglish / \setupRussian below — see end of preamble.

% --- ticket environment -------------------------------------------------------
% Renders the verbatim formulation from программа.pdf at the top of each topic.
% Visual styling only — language-aware label is unnecessary because each topic
% file is per-language already.
\newenvironment{ticket}{%
  \par\medskip\noindent\itshape%
}{%
  \par\medskip%
}

% --- bibliography (GOST 7.0.5/7.1) --------------------------------------------
\usepackage[
    backend=biber,
    style=gost-numeric,
    sorting=none,
    language=auto,
    autolang=other,
]{biblatex}
% Resolved relative to the compilation cwd. The build script (scripts/build.py)
% adds the repo root to BIBINPUTS, so this works for both book-* (cwd = root)
% and chapter-* (cwd = chapter dir, subfile compilation).
\addbibresource{literature.bib}

% --- hyperlinks & cross-refs --------------------------------------------------
\usepackage{hyperref}
\hypersetup{
    colorlinks=true,
    linkcolor=blue!50!black,
    citecolor=green!50!black,
    urlcolor=blue!70!black,
}
\usepackage[capitalize,nameinlink]{cleveref}

% --- Language setup commands -------------------------------------------------
% Each book file calls one of these immediately inside \begin{document}.
% They bundle babel's language switch with proof- and cleveref-name
% localization, so the result is the same regardless of which language
% babel happened to pick as "main".
\newcommand{\setupRussian}{%
  \selectlanguage{russian}%
  \renewcommand{\proofname}{Доказательство}%
  \crefname{theorem}{теорема}{теоремы}\Crefname{theorem}{Теорема}{Теоремы}%
  \crefname{lemma}{лемма}{леммы}\Crefname{lemma}{Лемма}{Леммы}%
  \crefname{proposition}{утверждение}{утверждения}\Crefname{proposition}{Утверждение}{Утверждения}%
  \crefname{corollary}{следствие}{следствия}\Crefname{corollary}{Следствие}{Следствия}%
  \crefname{definition}{определение}{определения}\Crefname{definition}{Определение}{Определения}%
  \crefname{example}{пример}{примеры}\Crefname{example}{Пример}{Примеры}%
  \crefname{problem}{задача}{задачи}\Crefname{problem}{Задача}{Задачи}%
  \crefname{remark}{замечание}{замечания}\Crefname{remark}{Замечание}{Замечания}%
  \crefname{section}{раздел}{разделы}\Crefname{section}{Раздел}{Разделы}%
  \crefname{chapter}{глава}{главы}\Crefname{chapter}{Глава}{Главы}%
  \crefname{equation}{уравнение}{уравнения}\Crefname{equation}{Уравнение}{Уравнения}%
}
\newcommand{\setupEnglish}{%
  \selectlanguage{english}%
  \renewcommand{\proofname}{Proof}%
  \crefname{theorem}{theorem}{theorems}\Crefname{theorem}{Theorem}{Theorems}%
  \crefname{lemma}{lemma}{lemmas}\Crefname{lemma}{Lemma}{Lemmas}%
  \crefname{proposition}{proposition}{propositions}\Crefname{proposition}{Proposition}{Propositions}%
  \crefname{corollary}{corollary}{corollaries}\Crefname{corollary}{Corollary}{Corollaries}%
  \crefname{definition}{definition}{definitions}\Crefname{definition}{Definition}{Definitions}%
  \crefname{example}{example}{examples}\Crefname{example}{Example}{Examples}%
  \crefname{problem}{problem}{problems}\Crefname{problem}{Problem}{Problems}%
  \crefname{remark}{remark}{remarks}\Crefname{remark}{Remark}{Remarks}%
  \crefname{section}{section}{sections}\Crefname{section}{Section}{Sections}%
  \crefname{chapter}{chapter}{chapters}\Crefname{chapter}{Chapter}{Chapters}%
  \crefname{equation}{equation}{equations}\Crefname{equation}{Equation}{Equations}%
}

% --- subfiles: each chapter compiles as a standalone PDF too ------------------
\usepackage{subfiles}

% --- misc ---------------------------------------------------------------------
\usepackage{enumitem}
\usepackage{booktabs}
% Font expansion needs scalable Cyrillic fonts (the `cm-super` Debian package
% or `texlive-fonts-extra`). Disable it to stay portable; protrusion still
% works with the default Computer Modern bitmap fonts.
\usepackage[expansion=false,protrusion=true,final]{microtype}
\usepackage{tikz}
.
%
% Build: pdflatex + biber. See README.md and scripts/build.py.
% Author note: keep this file in English — it's tooling, not content.
% =============================================================================

% --- encoding & languages -----------------------------------------------------
\usepackage[utf8]{inputenc}
\usepackage[T2A,T1]{fontenc}
% Both languages loaded; babel makes the LAST one the main language
% (here Russian) so cyrillic fontencoding bindings biblatex-gost relies
% on are set up at \begin{document}. Each book file then issues
% \selectlanguage{...} immediately inside \begin{document} to set its
% own active language for captions, hyphenation, dates, etc.
\usepackage[english,russian]{babel}

% --- page geometry ------------------------------------------------------------
\usepackage[a4paper, margin=2.5cm]{geometry}

% --- math packages ------------------------------------------------------------
\usepackage{amsmath, amssymb, amsthm, mathtools}
\usepackage{mathrsfs}     % \mathscr
\usepackage{bm}           % bold math

% --- math "house style" macros ------------------------------------------------
% Some of these short names (notably \C from T1 fontenc, and \Prob from
% other packages) are already defined when the preamble runs, so we clear
% them first before redefining. If you need the original meaning of \C
% (caron accent) elsewhere, use \v{c} instead.
\let\R\undefined
\let\C\undefined
\let\N\undefined
\let\Z\undefined
\let\Q\undefined
\let\F\undefined
\let\E\undefined
\let\Prob\undefined

% Number sets.
\newcommand{\R}{\mathbb{R}}
\newcommand{\C}{\mathbb{C}}
\newcommand{\N}{\mathbb{N}}
\newcommand{\Z}{\mathbb{Z}}
\newcommand{\Q}{\mathbb{Q}}
\newcommand{\F}{\mathbb{F}}

% Probability symbols.
\newcommand{\E}{\mathbb{E}}
\newcommand{\Prob}{\mathbb{P}}

% Operators (rendered upright).
\DeclareMathOperator{\rank}{rank}
\DeclareMathOperator{\im}{Im}
\DeclareMathOperator{\spn}{span}
\DeclareMathOperator{\diag}{diag}
\DeclareMathOperator{\tr}{tr}
\DeclareMathOperator{\sgn}{sgn}
\DeclareMathOperator{\Var}{Var}
\DeclareMathOperator{\Cov}{Cov}
\DeclareMathOperator*{\argmin}{arg\,min}
\DeclareMathOperator*{\argmax}{arg\,max}

% Vectors and matrices.
\newcommand{\vect}[1]{\mathbf{#1}}
\newcommand{\mat}[1]{\mathbf{#1}}
\newcommand{\T}{^{\mathsf{T}}}     % transpose

% --- theorem environments -----------------------------------------------------
% Names are localized via babel's \iflanguage in the headers below.
\theoremstyle{plain}
\newtheorem{theorem}{\iflanguage{russian}{Теорема}{Theorem}}[chapter]
\newtheorem{lemma}[theorem]{\iflanguage{russian}{Лемма}{Lemma}}
\newtheorem{proposition}[theorem]{\iflanguage{russian}{Утверждение}{Proposition}}
\newtheorem{corollary}[theorem]{\iflanguage{russian}{Следствие}{Corollary}}

\theoremstyle{definition}
\newtheorem{definition}[theorem]{\iflanguage{russian}{Определение}{Definition}}
\newtheorem{example}[theorem]{\iflanguage{russian}{Пример}{Example}}
\newtheorem{problem}[theorem]{\iflanguage{russian}{Задача}{Problem}}

\theoremstyle{remark}
\newtheorem*{remark}{\iflanguage{russian}{Замечание}{Remark}}

% Localized strings (proofname, cleveref names) are bundled into
% \setupEnglish / \setupRussian below — see end of preamble.

% --- ticket environment -------------------------------------------------------
% Renders the verbatim formulation from программа.pdf at the top of each topic.
% Visual styling only — language-aware label is unnecessary because each topic
% file is per-language already.
\newenvironment{ticket}{%
  \par\medskip\noindent\itshape%
}{%
  \par\medskip%
}

% --- bibliography (GOST 7.0.5/7.1) --------------------------------------------
\usepackage[
    backend=biber,
    style=gost-numeric,
    sorting=none,
    language=auto,
    autolang=other,
]{biblatex}
% Resolved relative to the compilation cwd. The build script (scripts/build.py)
% adds the repo root to BIBINPUTS, so this works for both book-* (cwd = root)
% and chapter-* (cwd = chapter dir, subfile compilation).
\addbibresource{literature.bib}

% --- hyperlinks & cross-refs --------------------------------------------------
\usepackage{hyperref}
\hypersetup{
    colorlinks=true,
    linkcolor=blue!50!black,
    citecolor=green!50!black,
    urlcolor=blue!70!black,
}
\usepackage[capitalize,nameinlink]{cleveref}

% --- Language setup commands -------------------------------------------------
% Each book file calls one of these immediately inside \begin{document}.
% They bundle babel's language switch with proof- and cleveref-name
% localization, so the result is the same regardless of which language
% babel happened to pick as "main".
\newcommand{\setupRussian}{%
  \selectlanguage{russian}%
  \renewcommand{\proofname}{Доказательство}%
  \crefname{theorem}{теорема}{теоремы}\Crefname{theorem}{Теорема}{Теоремы}%
  \crefname{lemma}{лемма}{леммы}\Crefname{lemma}{Лемма}{Леммы}%
  \crefname{proposition}{утверждение}{утверждения}\Crefname{proposition}{Утверждение}{Утверждения}%
  \crefname{corollary}{следствие}{следствия}\Crefname{corollary}{Следствие}{Следствия}%
  \crefname{definition}{определение}{определения}\Crefname{definition}{Определение}{Определения}%
  \crefname{example}{пример}{примеры}\Crefname{example}{Пример}{Примеры}%
  \crefname{problem}{задача}{задачи}\Crefname{problem}{Задача}{Задачи}%
  \crefname{remark}{замечание}{замечания}\Crefname{remark}{Замечание}{Замечания}%
  \crefname{section}{раздел}{разделы}\Crefname{section}{Раздел}{Разделы}%
  \crefname{chapter}{глава}{главы}\Crefname{chapter}{Глава}{Главы}%
  \crefname{equation}{уравнение}{уравнения}\Crefname{equation}{Уравнение}{Уравнения}%
}
\newcommand{\setupEnglish}{%
  \selectlanguage{english}%
  \renewcommand{\proofname}{Proof}%
  \crefname{theorem}{theorem}{theorems}\Crefname{theorem}{Theorem}{Theorems}%
  \crefname{lemma}{lemma}{lemmas}\Crefname{lemma}{Lemma}{Lemmas}%
  \crefname{proposition}{proposition}{propositions}\Crefname{proposition}{Proposition}{Propositions}%
  \crefname{corollary}{corollary}{corollaries}\Crefname{corollary}{Corollary}{Corollaries}%
  \crefname{definition}{definition}{definitions}\Crefname{definition}{Definition}{Definitions}%
  \crefname{example}{example}{examples}\Crefname{example}{Example}{Examples}%
  \crefname{problem}{problem}{problems}\Crefname{problem}{Problem}{Problems}%
  \crefname{remark}{remark}{remarks}\Crefname{remark}{Remark}{Remarks}%
  \crefname{section}{section}{sections}\Crefname{section}{Section}{Sections}%
  \crefname{chapter}{chapter}{chapters}\Crefname{chapter}{Chapter}{Chapters}%
  \crefname{equation}{equation}{equations}\Crefname{equation}{Equation}{Equations}%
}

% --- subfiles: each chapter compiles as a standalone PDF too ------------------
\usepackage{subfiles}

% --- misc ---------------------------------------------------------------------
\usepackage{enumitem}
\usepackage{booktabs}
% Font expansion needs scalable Cyrillic fonts (the `cm-super` Debian package
% or `texlive-fonts-extra`). Disable it to stay portable; protrusion still
% works with the default Computer Modern bitmap fonts.
\usepackage[expansion=false,protrusion=true,final]{microtype}
\usepackage{tikz}
